% common.tex -- shared preamble for all layout specimens.
% Each specimen sets its own geometry BEFORE % common.tex -- shared preamble for all layout specimens.
% Each specimen sets its own geometry BEFORE % common.tex -- shared preamble for all layout specimens.
% Each specimen sets its own geometry BEFORE % common.tex -- shared preamble for all layout specimens.
% Each specimen sets its own geometry BEFORE \input{common}.
% Engine: LuaLaTeX (needs fontspec for EB Garamond).

\usepackage{fontspec}
\setmainfont{EB Garamond}[Numbers={OldStyle,Proportional}]
\usepackage{microtype}
\setlength\emergencystretch{1em}
\usepackage{xcolor}
\definecolor{grapeink}{HTML}{5B2A4E}   % one accent color, used sparingly

% Three levels of notes, provided by bigfoot:
%   level 1: arabic numbers   (\footnote)
%   level 2: lowercase letters (\footnoteB)  -- notes on notes
%   level 3: symbols, run-in   (\footnoteC)  -- notes on notes on notes
% bigfoot stacks the series at the page foot in declaration order and
% lets a note of one series contain notes of the next.
\usepackage{bigfoot}
\DeclareNewFootnote{default}
\DeclareNewFootnote{B}[alph]
\DeclareNewFootnote[para]{C}[fnsymbol]
% Letters and symbols restart on every page, so they never run past z or
% past the nine fnsymbol marks.
\usepackage{perpage}
\MakePerPage{footnoteB}
\MakePerPage{footnoteC}

% Semantic note commands used by passage.tex. Each specimen maps them to
% a concrete layout, so the text is written once.
\providecommand\note[1]{\footnote{#1}}
\providecommand\subnote[1]{\footnoteB{#1}}
\providecommand\subsubnote[1]{\footnoteC{#1}}

% A plain, McPhee-like chapter opening: number, title, no ornament.
\newcommand\chapteropening[2]{%
  \clearpage\thispagestyle{empty}%
  \vspace*{0.18\textheight}%
  \begin{center}
    {\large\scshape\color{grapeink} #1}\par\medskip
    {\LARGE #2}\par
  \end{center}
  \vspace{2\baselineskip}\noindent\ignorespaces}

% Specimen label printed in the running foot, so the reader of the
% combined PDF knows which layout a page belongs to.
\newcommand\specimenlabel{}
\usepackage{fancyhdr}
\pagestyle{fancy}\fancyhf{}
\renewcommand\headrulewidth{0pt}
\fancyfoot[C]{\footnotesize\scshape\color{grapeink}\specimenlabel}
\fancypagestyle{empty}{\fancyhf{}\fancyfoot[C]{\footnotesize\scshape\color{grapeink}\specimenlabel}}
.
% Engine: LuaLaTeX (needs fontspec for EB Garamond).

\usepackage{fontspec}
\setmainfont{EB Garamond}[Numbers={OldStyle,Proportional}]
\usepackage{microtype}
\setlength\emergencystretch{1em}
\usepackage{xcolor}
\definecolor{grapeink}{HTML}{5B2A4E}   % one accent color, used sparingly

% Three levels of notes, provided by bigfoot:
%   level 1: arabic numbers   (\footnote)
%   level 2: lowercase letters (\footnoteB)  -- notes on notes
%   level 3: symbols, run-in   (\footnoteC)  -- notes on notes on notes
% bigfoot stacks the series at the page foot in declaration order and
% lets a note of one series contain notes of the next.
\usepackage{bigfoot}
\DeclareNewFootnote{default}
\DeclareNewFootnote{B}[alph]
\DeclareNewFootnote[para]{C}[fnsymbol]
% Letters and symbols restart on every page, so they never run past z or
% past the nine fnsymbol marks.
\usepackage{perpage}
\MakePerPage{footnoteB}
\MakePerPage{footnoteC}

% Semantic note commands used by passage.tex. Each specimen maps them to
% a concrete layout, so the text is written once.
\providecommand\note[1]{\footnote{#1}}
\providecommand\subnote[1]{\footnoteB{#1}}
\providecommand\subsubnote[1]{\footnoteC{#1}}

% A plain, McPhee-like chapter opening: number, title, no ornament.
\newcommand\chapteropening[2]{%
  \clearpage\thispagestyle{empty}%
  \vspace*{0.18\textheight}%
  \begin{center}
    {\large\scshape\color{grapeink} #1}\par\medskip
    {\LARGE #2}\par
  \end{center}
  \vspace{2\baselineskip}\noindent\ignorespaces}

% Specimen label printed in the running foot, so the reader of the
% combined PDF knows which layout a page belongs to.
\newcommand\specimenlabel{}
\usepackage{fancyhdr}
\pagestyle{fancy}\fancyhf{}
\renewcommand\headrulewidth{0pt}
\fancyfoot[C]{\footnotesize\scshape\color{grapeink}\specimenlabel}
\fancypagestyle{empty}{\fancyhf{}\fancyfoot[C]{\footnotesize\scshape\color{grapeink}\specimenlabel}}
.
% Engine: LuaLaTeX (needs fontspec for EB Garamond).

\usepackage{fontspec}
\setmainfont{EB Garamond}[Numbers={OldStyle,Proportional}]
\usepackage{microtype}
\setlength\emergencystretch{1em}
\usepackage{xcolor}
\definecolor{grapeink}{HTML}{5B2A4E}   % one accent color, used sparingly

% Three levels of notes, provided by bigfoot:
%   level 1: arabic numbers   (\footnote)
%   level 2: lowercase letters (\footnoteB)  -- notes on notes
%   level 3: symbols, run-in   (\footnoteC)  -- notes on notes on notes
% bigfoot stacks the series at the page foot in declaration order and
% lets a note of one series contain notes of the next.
\usepackage{bigfoot}
\DeclareNewFootnote{default}
\DeclareNewFootnote{B}[alph]
\DeclareNewFootnote[para]{C}[fnsymbol]
% Letters and symbols restart on every page, so they never run past z or
% past the nine fnsymbol marks.
\usepackage{perpage}
\MakePerPage{footnoteB}
\MakePerPage{footnoteC}

% Semantic note commands used by passage.tex. Each specimen maps them to
% a concrete layout, so the text is written once.
\providecommand\note[1]{\footnote{#1}}
\providecommand\subnote[1]{\footnoteB{#1}}
\providecommand\subsubnote[1]{\footnoteC{#1}}

% A plain, McPhee-like chapter opening: number, title, no ornament.
\newcommand\chapteropening[2]{%
  \clearpage\thispagestyle{empty}%
  \vspace*{0.18\textheight}%
  \begin{center}
    {\large\scshape\color{grapeink} #1}\par\medskip
    {\LARGE #2}\par
  \end{center}
  \vspace{2\baselineskip}\noindent\ignorespaces}

% Specimen label printed in the running foot, so the reader of the
% combined PDF knows which layout a page belongs to.
\newcommand\specimenlabel{}
\usepackage{fancyhdr}
\pagestyle{fancy}\fancyhf{}
\renewcommand\headrulewidth{0pt}
\fancyfoot[C]{\footnotesize\scshape\color{grapeink}\specimenlabel}
\fancypagestyle{empty}{\fancyhf{}\fancyfoot[C]{\footnotesize\scshape\color{grapeink}\specimenlabel}}
.
% Engine: LuaLaTeX (needs fontspec for EB Garamond).

\usepackage{fontspec}
\setmainfont{EB Garamond}[Numbers={OldStyle,Proportional}]
\usepackage{microtype}
\setlength\emergencystretch{1em}
\usepackage{xcolor}
\definecolor{grapeink}{HTML}{5B2A4E}   % one accent color, used sparingly

% Three levels of notes, provided by bigfoot:
%   level 1: arabic numbers   (\footnote)
%   level 2: lowercase letters (\footnoteB)  -- notes on notes
%   level 3: symbols, run-in   (\footnoteC)  -- notes on notes on notes
% bigfoot stacks the series at the page foot in declaration order and
% lets a note of one series contain notes of the next.
\usepackage{bigfoot}
\DeclareNewFootnote{default}
\DeclareNewFootnote{B}[alph]
\DeclareNewFootnote[para]{C}[fnsymbol]
% Letters and symbols restart on every page, so they never run past z or
% past the nine fnsymbol marks.
\usepackage{perpage}
\MakePerPage{footnoteB}
\MakePerPage{footnoteC}

% Semantic note commands used by passage.tex. Each specimen maps them to
% a concrete layout, so the text is written once.
\providecommand\note[1]{\footnote{#1}}
\providecommand\subnote[1]{\footnoteB{#1}}
\providecommand\subsubnote[1]{\footnoteC{#1}}

% A plain, McPhee-like chapter opening: number, title, no ornament.
\newcommand\chapteropening[2]{%
  \clearpage\thispagestyle{empty}%
  \vspace*{0.18\textheight}%
  \begin{center}
    {\large\scshape\color{grapeink} #1}\par\medskip
    {\LARGE #2}\par
  \end{center}
  \vspace{2\baselineskip}\noindent\ignorespaces}

% Specimen label printed in the running foot, so the reader of the
% combined PDF knows which layout a page belongs to.
\newcommand\specimenlabel{}
\usepackage{fancyhdr}
\pagestyle{fancy}\fancyhf{}
\renewcommand\headrulewidth{0pt}
\fancyfoot[C]{\footnotesize\scshape\color{grapeink}\specimenlabel}
\fancypagestyle{empty}{\fancyhf{}\fancyfoot[C]{\footnotesize\scshape\color{grapeink}\specimenlabel}}
